\chapter{Covariance, Invariance and Symmetry}
\label{ch:cis}
\epigraph{‘When I use a word,’ Humpty Dumpty said in rather a scornful tone,
‘it means just what I choose it to mean --- neither more nor less.’} {Lewis
Carroll}

Few words in physics are used as loosely as \emph{invariant}, \emph{covariant}
and \emph{symmetric}. One reads of the ``Lorentz invariance of Maxwell's
equations'' and, a page later, of their ``Lorentz covariance''; of ``general
covariance'' as the foundation of general relativity and, in the next
paragraph, as a triviality of notation; of ``gauge symmetry'' that is plainly
not a symmetry in the sense in which rotations are. The three notions are
tied together by one logical structure, but they live at \emph{different
levels of description}, and mixing the levels is a dependable source of
confusion.

\begin{itemize}[leftmargin=1.6em,itemsep=2pt]
\item \textbf{Invariance} is a property of an \emph{object}: a number, a
  function, a tensor field, a functional, a particular solution. The object is
  left unchanged by a transformation.
\item \textbf{Covariance} is a property of a \emph{formulation}: a family of
  quantities, or an equation built from them, transforms in a definite,
  prescribed way (a representation of the transformation group), so that the
  equation keeps its form.
\item \textbf{Symmetry} is a property of a \emph{law} (a dynamical system): a
  transformation that carries solutions into solutions, or leaves the action
  unchanged.
\end{itemize}

Symmetry is the only one of the three that is a claim about the real world. It
says that if an experiment on an isolated system is repeated after the
\emph{entire} arrangement has been displaced, rotated or boosted, the outcome
is the transformed outcome. That claim can be, and has been, tested with
extraordinary precision (Michelson--Morley-type experiments for boosts,
Hughes--Drever-type experiments for rotations). Covariance, by contrast, is
a statement about how we write equations; whether it carries any physical
content depends on what else is held fixed.

A \emph{passive} transformation relabels the points of spacetime,
$x\mapsto x'=f(x)$, and asks how the \emph{same} physical field is described
in the new coordinates. For a scalar and for a vector field,
\begin{equation}\label{passive}
  \phi'(x') = \phi(x),\qquad
  A'^\mu(x') = \frac{\partial x'^\mu}{\partial x^\nu}\,A^\nu(x).
\end{equation}
An \emph{active} transformation by $g$ moves the field itself, keeping the
coordinates: for $g\colon x\mapsto \Lambda x + a$,
\begin{equation}\label{active}
  (g\cdot\phi)(x) = \phi(g^{-1}x),\qquad
  (g\cdot A)^\mu(x) = \Lambda^\mu{}_\nu\,A^\nu(g^{-1}x).
\end{equation}
The new field \emph{function} is given by the same formula in both cases
(solve \eqref{passive} for $\phi'(y)=\phi(f^{-1}y)$ and compare with
\eqref{active}); what differs is the interpretation. In the passive
reading one compares two descriptions of one physical situation; in the
active reading one compares two different physical situations in one
description. Symmetry is an active notion: it asks whether the second
situation is also allowed by the law. Covariance is naturally passive: it asks
whether the \emph{description} keeps its form.


The active reading carries a warning that will turn out to be the heart of
the matter: \emph{everything that enters the equations must be transformed,
or else be invariant.}

\begin{example}[Coulomb field of a fixed charge]\label{ex:coulomb}
The equation $\nabla\!\cdot{\bm{E}}=4\pi\rho$ with a prescribed source
$\rho=q\,\delta^3({\bm{x}}-{\bm{x}}_0)$ has the solution
${\bm{E}}({\bm{x}})=q({\bm{x}}-{\bm{x}}_0)/|{\bm{x}}-{\bm{x}}_0|^3$. Translating the field alone,
${\bm{E}}_a({\bm{x}})={\bm{E}}({\bm{x}}-\bm a)$, gives a field whose divergence is
$4\pi q\,\delta^3({\bm{x}}-{\bm{x}}_0-\bm a)$: it is \emph{not} a solution for the
original source. If the source is translated too, the pair
$({\bm{E}}_a,\rho_a)$ is a solution. Translation is therefore a symmetry of the
closed system ``field $+$ charge,'' and is broken as soon as the charge is
treated as a fixed, external, non-dynamical object.
\end{example}

Whatever is held fixed and not transformed is called a \emph{background} (or
\emph{absolute}) structure: a prescribed source, a fixed metric, a preferred
velocity field. Prescribing backgrounds is unavoidable in practice, and it is
precisely what separates covariance from symmetry.

Let us see how it works in classical mechanics which is
one of the most important models in physics.
Until the end of the 19th century, mechanics was considered as direct and the only
the basis of all physics. By words understand or
to explain a physical phenomenon meant constructing a mechanical model of it.

Situation has changed,  other physical models exist, e.g. quantum mechanics so 
"usual" mechanics has changed its name to "classical mechanics" and still is
the central part of theoretical physics. All main physical concepts emerged from
it.

\section{Two Roles of Classical Mechanics}
\label{sec:cis-2r}
Classical mechanics occupies a peculiar position among physical theories. It is,
first, a complete and self-contained subject: a body of concepts, equations, and
results describing the motion of bodies under forces, valid and useful on its own
terms, with no reference whatsoever to the Planck constant $\hbar$ or the speed of
light $c$. Celestial mechanics, rigid-body dynamics, the theory of small
oscillations, continuum mechanics, and the modern theory of dynamical systems and
chaos are all classical mechanics pursued for its own sake, and none of them is
waiting to be corrected by a more fundamental theory in order to be meaningful.

Second,  classical mechanics is
the \emph{base} on which relativistic and quantum mechanics are built. Special and
general relativity are commonly introduced as corrections to Newtonian mechanics
that become significant when velocities approach $c$ or fields become strong;
quantum mechanics is commonly introduced as a modification that becomes significant
at the scale of $\hbar$. In both cases, the classical theory is treated as a known,
well-understood limit --- the solid ground from which the generalization departs
and to which it must reduce.

It is tempting to think that once this second, subordinate role is granted, nothing
further needs to be said about classical mechanics: its job is simply to sit still
and be a limit. Here we argue that this is false in two related ways.
\begin{enumerate}
\item Several phenomena that are routinely presented as characteristically
relativistic or characteristically quantum-mechanical --- an underdetermination
of physical ontology by the observable facts, an ambiguity in the passage from
classical description to quantized theory, a geometric ``phase'' with no
classical counterpart, an apparent conflict between microscopic reversibility
and macroscopic causal, dissipative behavior --- already occur inside classical
mechanics, before relativity or quantization ever enter the discussion. When
such a paradox survives the excision of everything relativistic or quantum, its
source cannot honestly be attributed to relativity or to quantum theory.
\item Even restricted to the narrow task of serving as the base for
relativistic and quantum generalization, classical mechanics is not a settled,
transparent starting point. Producing the structure that gets generalized --- a
well-defined phase space, a Hamiltonian, a consistent variational principle ---
is itself a nontrivial achievement that fails, or becomes ambiguous, for entire
classes of ordinary classical systems.
\end{enumerate}

This book is devoted to the second role of classical mechanics, some
interesting problems beyond this role are briefly discussed in Appendices.

Classical mechanics is usually presented in one of two logical orders:
\begin{description} \item[Newton first.] The postulates are equations of motion
            (together with a catalogue of force laws and a statement about
            constraint reactions). Lagrange's equations, the principle of least
            action and the Hamiltonian formalism are \emph{derived}. This is
        the order of Goldstein \cite{GoldsteinPooleSafko2000}, Arnold's \cite{Arnold1989}
    first chapters, and most introductory courses. 

\emph{Strengths.} Every step has a physical referent. Students see where
the action comes from and, crucially, \emph{when it does not exist}.
Constraint reactions are computed, not suppressed. The nonholonomic
and dissipative cases are included without a change of philosophy.

\emph{Weaknesses.} The passage to generalized coordinates is algebraically
heavy. The reader emerges thinking Lagrangian mechanics is a
\emph{computational device}, and misses that the action is a primary
object in field theory and quantum mechanics.

\item[Action first.] The
        postulate is Hamilton's principle, $\delta S=0$ with $S=\int L\,dt$.
        Newton's equations are \emph{derived}, and conservation laws follow
        from symmetries of $L$ via Noether's theorem. This is the order of
        Landau and Lifshitz \cite{LandauLifshitz-1}, Lanczos \cite{Lanczos1970}
and many others. 

\emph{Strengths.} Symmetry becomes the organizing principle:
Noether's theorem explains energy, momentum and angular momentum at once.
Coordinate independence is manifest: $\delta S=0$ does not know what
coordinates are. The road to field theory and quantization is short.

\emph{Weaknesses.} The phrase ``least action'' invites teleology. The student
learns that dissipative systems are ``not mechanics'' and that constraint
reactions are an afterthought. Most damagingly, the nonholonomic case produces
a \emph{wrong answer} that is indistinguishable, in form, from the right one.
\end{description} 

Both orders lead to the same equations for simple systems of point particles,
so it is tempting to regard the choice as a matter of taste. It is not. The two
roads are \emph{not} mutually inverse once we leave the simplest class of
systems: the road from Newton to the action requires hypotheses (potential
forces, ideal constraints) that the road back silently assumes, and the two
roads give \emph{different} answers for nonholonomic constraints. 

In the language of chapter \ref{rm}, the Newtonian laws are statements anchored
in the real world: forces are measured, constraints are realized by hinges,
rails and blades. The action is a construction of the ideal world: a functional
whose extremals are, by design, trajectories. Both routes have their strengths
and weaknesses.

In section~\ref{sec:cis-intro} I follow \cite{Arnold1989} as more physical, later
I shall use mainly second (action first) approach with some reservations.
I see Hamilton's principle, $\delta S=0$ not as law of nature but as 
powerful heuristic.

\section{Introduction}
\label{sec:cis-intro}
In Ideal world of classical mechanics space is three-dimensional  and euclidean
$\mathbb{R}^3$ and time is one-dimensional $\mathbb{T}$.
So universe of classical mechanics is $\mathbb{T} \times \mathbb{R}^3$. 

The simplest object of classical mechanics is point mass (particle) i.e. object
whose shape, sizes and internal structure are inessential, so mass (positive
scalar $m$) is the only parameter. 
The state of a particle is fully determined by position ($\bm x \in\mathbb{R}^3$) 
and velocity ($\bm v \in\mathbb{R}^3$). 
More accurately: A motion in $\mathbb{R}^3$ is a differentiable mapping 
\[\bm x: \mathbb{I}\rightarrow \mathbb{R}^3, \mathbb{I} \subset \mathbb{R}.\] 
The derivative 
\[\dot {\bm x}(t)=\frac{d{\bm x}}{dt}(t)\]
is called the velocity vector. 
The derivative 
\[\ddot {\bm x}(t)=\frac{ d^2 {\bm x}}{dt^2}(t)\]
is called the acceleration vector. 
The image of motion (curve in $\mathbb{R}^3$) called trajectory and  curve in 
$\mathbb{T} \times \mathbb{R}^3$ which appears  as the graph of a motion, is
called a world line.

The initial state of a mechanical system (the totality of positions and
velocities of its points at some moment of time) uniquely determines all of
its motion. In particular, the initial positions and velocities determine the
accelerations. This is called `Newton's principle of determinacy'.

Let us consider the mechanical system consisting of $n$ particles 
with masses $m_j,j=1\ldots n$ and word lines $\bm x_j(t)$.
Suppose that $n$ vector functions   
\[{\bm F}_j:\mathbb{R}^{3 n}\times\mathbb{R}^{3 n}
\times\mathbb{R}\rightarrow\mathbb{R}^3,j=1\ldots n\]
are defined. Then the motion of all particles is determined by equations
\begin{equation}
  \label {newton-eq}
    m_j \ddot {\bm x}_j={\bm F}_j({\bm x}_k,\dot {\bm x}_k,t) \quad j,k=1\ldots n 
\end{equation}
and initial conditions
\begin{equation}
\label {newton-ic}
    {\bm x}_j(t_0)={\bm y}_j\quad\dot {\bm x}_j(t_0)={\bm v}_j  \quad j=1\ldots n 
\end{equation}

Newton's equation \eqref{newton-eq} is  the basis of mechanics. Functions 
$\bm F_j$ are called forces they should be determined  for each specific 
mechanical system. If our system is isolated i.e. includes explicitly all 
interacting particles, then galilean invariance puts some restrictions 
on forces.

Namely 
\begin{equation*}
     m_j \ddot {\bm x}_j=
    {\bm F}_j({\bm x}_j-{\bm x}_k,\dot {\bm x}_j-\dot {\bm x}_k)\quad j,k=1\ldots n 
\end{equation*}
practically such systems occur in celestial mechanics. In most other cases the
system considered is not isolated and forces of general kind \eqref{newton-eq} 
appear as a result of interaction with environment.

Some forces can not be set explicitly they appear as a result of constraints
imposed on the system and should be figured out form the prescribed movement.
\begin{equation}
  \label {newton-eq-const}
    m_j \ddot {\bm x}_j={\bm F}_j({\bm x}_k,\dot {\bm x}_k,t) 
    +{\bm R}_j \quad j,k=1\ldots n 
\end{equation}
These constraint forces ${\bm R}_j$ are result of collective action of 
\emph{all constraints}. If a system consists of mass points connected by 
rods, hinges, etc., then the need may arise to talk about the constraint force 
of some particular constraint. However the concept of a constraint force for an 
individual constraint is generally \emph{impossible to define}.

Mathematics of  Cauchy problem (\ref{newton-eq},\ref{newton-ic}) 
one can find in e.g. \cite{Arnold1992}. Generally the unique solution exists if
all functions are smooth enough. However smoothness may be violated in 
physically important cases, anyway functions $\bm x_j(t)$ should be continuous,
$\bm v_j(t)$  may have finite jumps as a results of strong force acting for 
very short time (e.g. elastic collision).

Note that initial conditions have nothing to do with causality. Causality makes
sense in Real world only, in Ideal world we are free to set final conditions or
fix $\bm x_j(t_1)$ and $\bm v_j(t_1)$ for some intermediate time $t_1$. In
many cases we are interested in solution of equation \eqref{newton-eq} with
fixed initial velocities and final coordinates  (aiming) or fixed coordinates
on both ends (boundary problem). Anyway when conditions are set at two
different times  solution may be not unique (unless force is changing
relatively slow) \footnote {I discovered this fact around 1980 and was very
proud of myself until I found papers \cite{Maslov1962a} and \cite{Maslov1962b}.}.

If we allow only potential forces i.e. 
\begin{equation}
  \label {pot-force}
        \bm F_j(\bm x_k,t)=-\frac{\partial U(\bm x_k,t)}{\partial \bm x_j} \quad j,k=1\ldots n 
\end{equation}
we shall see that equations of motion (\ref{newton-eq}) can be written as 
\begin{equation}
   \label {lagr-eq}
    \frac{d}{dt} \frac{\partial L({\bm x}_k,{\bm v}_k,t)}{\partial {\bm v}_j}= 
    \frac{\partial L({\bm x}_k,{\bm v}_k,t)}{\partial {\bm x}_j}\quad j,k=1\ldots n 
\end{equation} 
where 
\begin{equation}
  \label{lagr-nat}
        L({\bm x}_k,{\bm v}_k,t)=\Sigma_1^n \frac{m_k}{2} {\bm v}_k \cdot {\bm v}_k-U({\bm x}_k,t) 
\end{equation}
Function $L({\bm x}_k,{\bm v}_k,t)$ is called Lagrangian, vectors 
\[{\bm p}_j= \frac{\partial L({\bm x}_k,{\bm v}_k,t)}{\partial {\bm v}_j}
=m_j \bm v_j \] 
are called momenta. 
We can write the equations \eqref{lagr-eq} in more symmetrical way:
\begin{equation}
   \label {ham-eq}
   \begin{split}
       \dot {\bm x}_j&=\frac{\partial H(\bm x_k,\bm p_k,t)}{\partial \bm p_j}  \\ 
        \dot {\bm p}_j&= -\frac{\partial H(\bm x_k,\bm p_k,t)}{\partial \bm x_j}\quad j,k=1\ldots n    
   \end{split}
\end{equation}
where 
\begin{equation}
        H(\bm x_k,\bm p_k,t)=\Sigma_1^n \frac{1}{2 m_k} \bm p_k \cdot \bm p_k+U(\bm x_k,t) \label{ham-nat}
\end{equation}
Note that functions $L$ (Lagrangian) and $H$ (Hamiltonian) are connected by Legendre transform:
\begin{equation}
        L(\bm x_k,\bm v_k,t)=\Sigma_1^n  \bm p_k \cdot \bm v_k -H(\bm x_k,\bm p_k,t)|_{\bm p_k=m_k\bm v_k}
\quad k=1\ldots n \label {lagr-ham}
\end{equation}


\begin{definition}
        The equation
        \[\frac{d}{d t}(\frac{\partial L}{\partial \dot {\bm x}})-\frac{\partial L}{\partial  {\bm x}}=0 \]
        is called the Euler-Lagrange equation for the functional
        \[S=\int_{t_1}^{t_2} L(\bm x,\dot {\bm x},t) dt \]
\end{definition}
Let \( \gamma= \bm x(t) \quad t_1 \le t \le t_2 \) be a curve  in  $\mathbb{R} \times \mathbb{R}^n$ and 
$L: \mathbb{R}^n \times \mathbb{R}^n \times \mathbb{R}\rightarrow \mathbb{R}$ a function of $2 n +1$
variables, then
\begin{theorem}
    The curve $\gamma$ is an extremal of the functional \[S=\int_{t_1}^{t_2} L(\bm x,\dot {\bm x},t) dt \]
        on the space of curves joining $(t_1, \bm x_1)$ and $(t_2,\bm x_2)$, if and only if the Euler-
Lagrange equation is satisfied along $\gamma$.
\end{theorem}
This is a system of n second-order equations, and the solution depends on
$2 n$ arbitrary constants. The $2 n$ conditions $\bm x(t_1) = \bm x_0$ and $\bm x(t_2) = \bm x_2$ are used
for finding them.
Two important remarks:
\begin{enumerate}
        \item There may be many extremals connecting two
given points or none at all.
\item The condition for a curve  to be an extremal of a functional does not depend
on the choice of coordinate system. This paves way to geometrisation of mechanics.
\end{enumerate}

It is easy to see that equation of motion (\ref{lagr-eq}) are actually
Euler-Lagrange equations for lagrangian \eqref{lagr-nat}. So we have  Hamilton's
principle of least action.

Equations of motion (EOM) are the heart of classical mechanics in both its roles.
However equations themselves are not enough, to get meaningful solution
one should fix the \emph{problem} i.e. specify the conditions imposed.
The problems usually  associated with classical mechanics are:
\begin{description}
    \item [IVP / Cauchy problem]  The full state (position and velocity, or a
  phase-space point) at one instant $t_0$ is specified.
   The Picard--Lindel\"of theorem gives
  local existence and uniqueness for Lipschitz forces; global existence
  requires a no-blow-up bound.
   This is what time-stepping integrators  actually solve.
Two large class of methods are in use: 
\begin{enumerate} 
    \item "Usual" one-time integrators when differential equations are replaced
        with finite difference ones defining snapshots at discrete time steps.
        These integrators allow accurate description of standard problem:
        initial conditions, finite time interval. 
\item Geometric integrators which are less accurate for any single trajectory
    but accurately present qualitative picture for large set of initial data
        and long times. Such integrators may calculate  coordinates and momenta
        at different times
\end{enumerate}
One can look for details in \cite{HairerLubichWanner2006} for example.

    \item [Boundary Value Problem (BVP)] Data at two or more distinct times/points,
  e.g. $x(t_0) = x_0$, $x(t_1) = x_1$ --- not a full state at one
  instant is fixed.
   Existence and uniqueness are \emph{not} automatic (unlike the
  IVP): zero, one, or many solutions can occur.
   \textbf{Prototypes:} the brachistochrone problem, geodesics
  between two points, Lambert's problem, and any Euler--Lagrange
  equation arising from a fixed-endpoint action principle.
   \textbf{Numerically:} shooting methods, relaxation /
  finite-difference methods, collocation.

    \item [Eigenvalue Problem (a BVP subtype)] A linear EOM together with boundary conditions admits
  nontrivial solutions only for special parameter values ---
  normal-mode frequencies, quantized energies.

    \item [Periodicity Condition] Search for $x(t+T) = x(t)$. Effectively a BVP with the two
  endpoints identified, but $T$ is itself unknown, making it a
  nonlinear eigenvalue-type problem.
   \textbf{Poincar\'e--Bendixson theorem} (2D flows): a bounded
  trajectory that does not approach a fixed point must approach a
  periodic orbit.

    \item [Control / Optimization Problem] An external input, or a member of a family of
  EOMs, that steers the system to a target state while extremizing a
  cost functional --- Pontryagin's maximum principle is sought.
   Structurally a BVP with an added adjoint (costate) system.

    \item [Scattering Problem] An asymptotic free state as $t \to
  -\infty$ (or $r \to \infty$) is specified.
   The asymptotic outgoing state as $t \to
  +\infty$ --- the $S$-matrix or cross-section --- given a localized
  interaction region in between is sought.
   This is asymptotic BVP: the conditions sit at $t = \pm\infty$ rather
   than at finite times, which   is exactly why it needs its own machinery 
   ($S$-matrix, asymptotic
  completeness) rather than shooting/relaxation methods.
 Often the trajectory itself is discarded --- only the
  in $\to$ out map matters.

    \item [Inverse Problem] Having observed trajectories or scattering data,
   one searches  the force law / potential / EOM itself
  (inverse scattering; Newton inferring $1/r^2$ from Kepler's laws is
  the archetype).
\end{description}

\section{Finite and infinite numbers of degrees of freedom}
\label{sec:dof}
Usually mechanical systems (finite number of degrees of freedom) 
and field theories (infinite number of degrees of freedom) are seen
as absolutely different. However form symmetry (or covariance) point
of view the two are
not different theories: a mechanical system is a field theory on a
one-dimensional base manifold, and a field theory is a mechanical system whose
coordinates carry a continuous label. This section makes both statements
explicit and notes where the reduction fails to be a pure special case.

\paragraph{Mechanics as a field theory on $M=\mathbb R_t$.}
Take the base manifold to be the time axis, with the single coordinate
$x^0\equiv t$, and identify
\begin{equation}\label{dictionary-id}
  \varphi_a(x^\mu)\;\leftrightarrow\;q^i(t),\qquad
  \mathcal L(\varphi_a,\partial_\mu\varphi_a;x)\;\leftrightarrow\;L(q^i,\dot q^i,t),\qquad
  \partial_\mu\;\leftrightarrow\;\frac{\dd}{\dd t}.
\end{equation}
The formulas of Sections~\ref{sec:definitions} and~\ref{sec:noether} then
specialize as in Table~\ref{tab:dictionary-dof}.

\begin{table}[ht]
\centering
\small
\renewcommand{\arraystretch}{1.35}
\setlength{\tabcolsep}{4pt}
\begin{tabular}{@{}>{\raggedright\arraybackslash}p{7.3cm}>{\raggedright\arraybackslash}p{7.3cm}@{}}
\toprule
Field theory on $M^d$ & Mechanics ($d=1$)\\
\midrule
Fields $\varphi_a(x^\mu)$ & Coordinates $q^i(t)$\\
Euler--Lagrange: $\pd{\mathcal L}{\varphi_a}-\partial_\mu\pd{\mathcal L}{(\partial_\mu\varphi_a)}=0$
  & $\pd{L}{q^i}-\dfrac{\dd}{\dd t}\pd{L}{\dot q^i}=0$\\
Noether current: $J^\mu=\pd{\mathcal L}{(\partial_\mu\varphi_a)}X_a-K^\mu$,\ \ $\partial_\mu J^\mu\approx0$
  & Noether charge: $Q=p_iX^i-K$,\ \ $\dot Q\approx0$\\
Energy--momentum tensor: $T^\mu{}_\nu=\pd{\mathcal L}{(\partial_\mu\varphi_a)}\partial_\nu\varphi_a-\delta^\mu_\nu\mathcal L$
  & $T^0{}_0=p_i\dot q^i-L=H$\\
Explicit $x$-dependence: $\partial_\mu T^\mu{}_\nu\approx-\big(\pd{\mathcal L}{x^\nu}\big)_{\!\rm expl}$
  & $\dot H=-\pd{L}{t}$\\
\bottomrule
\end{tabular}
\caption{Mechanics as the case $d=1$ of field theory.}
\label{tab:dictionary-dof}
\end{table}

In $d=1$ the divergence $\partial_\mu J^\mu$ is a single derivative, so the
current \emph{is} the charge and the conservation law is constancy in time.
The derivation of Noether's first theorem is the same line in both cases. For
$\delta q^i=\eps X^i$ with $\delta L=\eps\dot K$,
\begin{equation}\label{noether-mech}
  \delta L=\Big[\pd{L}{q^i}-\frac{\dd}{\dd t}\pd{L}{\dot q^i}\Big]\eps X^i
  +\eps\,\frac{\dd}{\dd t}\big(p_iX^i\big)=\eps\dot K
  \;\Longrightarrow\;
  \frac{\dd}{\dd t}\big(p_iX^i-K\big)\approx0,\qquad p_i=\pd{L}{\dot q^i},
\end{equation}
which is \eqref{noether1} with one base coordinate. The Galilean boost of
Section~\ref{sec:galilei}, $X^i=u^it$ and $K=m\,\bm u\!\cdot\!{\bm{x}}$, is an
application of exactly this formula.

\paragraph{The covariance machinery in $d=1$.}
Proposition~\ref{prop:main} applies with $\mathcal G=\diff(\mathbb R)$, the
reparametrizations of time, and with the preferred chart, that is the one-form
$\dd t$, as the background $b_0$. A reparametrization $t\to f(t)$ pulls $\dd t$
back to $f'(t)\,\dd t$, so
\begin{equation}
  \stab_{\diff(\mathbb R)}(\dd t)=\{\,t\to t+c\,\}.
\end{equation}
The symmetry is time translation, and Noether's theorem gives conservation of
$H$. A prescribed time-dependent external force is a second background; it
shrinks the stabilizer to the identity and produces $\dot H=-\partial L/\partial t$,
the last row of Table~\ref{tab:dictionary-dof}. Noether's \emph{second}
theorem applies to the parametrized action of Example~\ref{ex:param} (or of
Section~\ref{sec:sr}), where $t$ is promoted to a dynamical variable: the
identity \eqref{param-identity} is \eqref{noether2} for
$\diff(\mathbb R)$. The soldering form of Section~\ref{sec:soldering} reduces
to the single function $e_\lambda=\dd X^0/\dd\lambda=\dot t$, with
$D=\dot t^{-1}\dd/\dd\lambda$; the parametrized Newtonian action is
$\int\dot t\,\tfrac m2(D{\bm{x}})^2\,\dd\lambda$, which equals $\int\tfrac m2\dot{\bm{x}}^2/\dot t\,\dd\lambda$.

\paragraph{Base symmetries and target symmetries.}
A mechanical system is a map $q\colon M^1\to N$ from the time axis into a
\emph{target} space $N$ (a sigma model). Transformations then split into those
acting on the base, which change the arguments of the fields, and those acting
on the target, which change their values. Table~\ref{tab:base-target} shows
that the distinction covers every example of the chapter.

\begin{table}[ht]
\centering
\footnotesize
\renewcommand{\arraystretch}{1.3}
\setlength{\tabcolsep}{3pt}
\begin{tabular}{@{}>{\raggedright\arraybackslash}p{3.0cm}%
                  >{\raggedright\arraybackslash}p{2.1cm}%
                  >{\raggedright\arraybackslash}p{2.4cm}%
                  >{\raggedright\arraybackslash}p{3.0cm}%
                  >{\raggedright\arraybackslash}p{3.6cm}@{}}
\toprule
System & Base $M$ & Target / fibre & Base symmetries & Target symmetries\\
\midrule
Newtonian particle & $\mathbb R_t$ & $E^3$ & time translation
  & $E(3)$; boosts are target maps with $t$-dependent parameter\\
Relativistic particle & $\mathbb R_\lambda$ & Minkowski space
  & $\diff(\mathbb R)$ (gauge) & Poincar\'e\\
Real scalar field & Minkowski space & $\mathbb R$ & Poincar\'e
  & $\varphi\to-\varphi$ if $V$ is even\\
$O(N)$ scalar model & Minkowski space & $\mathbb R^N$ & Poincar\'e & $O(N)$\\
Gravity (tetrad form) & $M^4$ & fields $e^a{}_\mu$ & $\diff(M)$ (gauge)
  & local Lorentz (gauge)\\
\bottomrule
\end{tabular}
\caption{Base and target symmetries.}
\label{tab:base-target}
\end{table}

\paragraph{Field theory as a mechanical system with a continuum of degrees of
freedom.}
In the opposite direction, a $3{+}1$ split of a field theory turns it into a
mechanical system: set $q_{{\bm{x}}}(t)=\varphi({\bm{x}},t)$, so that the discrete index
$i$ becomes the continuous label ${\bm{x}}$. Sums become integrals,
\begin{equation}
  L=\int\dd^3x\,\mathcal L,\qquad
  \pi({\bm{x}})=\frac{\delta L}{\delta\dot\varphi({\bm{x}})},\qquad
  Q=\sum_ip_iX^i-K\ \longrightarrow\ \int\dd^3x\,\big(\pi X-K^0\big)=\int\dd^3x\,J^0 .
\end{equation}
A chain of masses shows how the two descriptions relate. On a ring of $N$
masses,
\begin{equation}
  L=\sum_{n}\Big[\tfrac12m\,\dot q_n^2-\tfrac12k\,(q_{n+1}-q_n)^2\Big].
\end{equation}
Set $q_n(t)=\varphi(na,t)$, $m=\rho a$, $k=\tau/a$, and keep the length $Na$
fixed. As $a\to0$,
\begin{equation}
  L\ \to\ \int\dd x\,\Big[\tfrac12\rho\,\dot\varphi^2-\tfrac12\tau\,\varphi'^2\Big],
  \qquad\rho\ddot\varphi=\tau\varphi'' .
\end{equation}
The same system has three different symmetry statements.
\begin{enumerate}[leftmargin=1.8em,itemsep=3pt,label=\textbf{\arabic*.}]
\item \emph{Shift} $q_n\to q_n+\eps$ for all $n$, a target symmetry with $X=1$,
  $K=0$. The charge $\sum_np_n=\sum_nm\dot q_n$ becomes $\int\rho\dot\varphi\,\dd x$,
  with the local current $J^t=\rho\dot\varphi$, $J^x=-\tau\varphi'$, and
  $\partial_tJ^t+\partial_xJ^x=\rho\ddot\varphi-\tau\varphi''=0$.
\item \emph{Lattice translation} $n\to n+1$, a discrete symmetry (the cyclic group
  $\mathbb Z_N$ of the ring). It has no Noether charge in the sense of
  Section~\ref{sec:noether}.
\item \emph{Continuum translation} $x\to x+b$, a base symmetry that exists only
  in the limit. By \eqref{noether1} with $\delta\varphi=-b\,\partial_x\varphi$
  its charge is the field momentum $P=-\int\rho\,\dot\varphi\,\varphi'\,\dd x$.
\end{enumerate}
So the relation between the finite and the infinite system is not only
``special case.'' The finite chain has \emph{fewer} symmetries than its
continuum limit, because discretization generically breaks continuous base
symmetries (and, for a relativistic theory on a lattice, Lorentz invariance).

\paragraph{Where the reduction is not a pure special case.}
\begin{itemize}[leftmargin=1.6em,itemsep=2pt]
\item \emph{Boundary conditions.} For $d=1$ the only boundaries are the initial
  and final times, and $Q$ needs no fall-off condition. For $d>1$ the charge
  $Q=\int J^0\,\dd^3x$ is conserved only if the flux through spatial infinity
  vanishes, which requires conditions on the fields.
\item \emph{Constraints.} A gauge parameter $\eps({\bm{x}},t)$ yields one constraint
  per spatial point (Gauss's law, $\nabla\!\cdot{\bm{E}}=4\pi\rho$, in
  Example~\ref{ex:noether2-em}), whereas reparametrization in $d=1$ yields a
  single constraint.
\item \emph{Soldering form.} It has $d$ components ($d=4$ for spacetime), not one.
\item \emph{Manifest covariance.} A $3{+}1$ split gives up manifest covariance,
  and boosts then act nonlinearly, exactly as for the gauge-fixed particle
  $\lambda=t$ of Section~\ref{sec:sr}.
\item \emph{Symmetry breaking in the quantum theory.} With finitely many
  degrees of freedom the Stone--von Neumann theorem makes all irreducible
  representations of the canonical commutation relations unitarily
  equivalent, so a symmetry is implemented unitarily and a degenerate classical
  minimum is resolved by tunneling (the ground state of a symmetric double well
  is symmetric). With infinitely many degrees of freedom inequivalent
  representations exist, and spontaneous symmetry breaking of the quantum
  vacuum (Section~\ref{sec:kepler}) becomes possible. Classically, the
  distinction between symmetry of the law and symmetry of the solution
  (Section~\ref{sec:law-solution}) is already present in both cases.
\end{itemize}
In both directions Proposition~\ref{prop:main} is unchanged: the symmetry group
is the stabilizer of the background in the covariance group. Only the
dimension of the base and the supply of target symmetries differ.

\section{Definitions and the basic theorem}
\label{sec:definitions}

A group $G$ acts on a vector space $V$ through a \emph{representation},
a map $D\colon G\to \mathrm{GL}(V)$ with $D(g_1g_2)=D(g_1)D(g_2)$.
Fields are classified by the representation of the spacetime group under
which they transform: scalars (trivial), vectors, tensors, spinors. A scalar
field transforms by the first formula of \eqref{active}, a vector field by
the second, with $D(g)=\Lambda$. The trivial representation $D(g)=1$ is
special: objects transforming in it are \emph{invariant}. Invariance is thus
the limiting case of covariance in which the representation is trivial.

Let $\phi$ denote the dynamical fields (collectively) and $b$ the background
fields, both defined on spacetime $M$ (for a mechanical system $M$ is the time axis;
see Section~\ref{sec:dof}). A group $\mathcal G$ acts on pairs,
$g\colon(\phi,b)\mapsto(g\cdot\phi,\,g\cdot b)$. A system of field equations
is an operator $E$,
\begin{equation}\label{EOM}
  E[\phi;\,b]=0 .
\end{equation}

\begin{definition}[Invariance]\label{def:invariance}
An object $T$ (a field, a functional, a tensor, a solution) is
\emph{invariant} under a group $G$ if $g\cdot T=T$ for all $g\in G$.
\end{definition}

Familiar examples: the metric $\eta_{\mu\nu}$ is invariant under the
Lorentz group, which is exactly the content of \eqref{poincare}; $\delta_{ij}$
and $\epsilon_{ijk}$ are invariant under $SO(3)$; the interval
$\eta_{\mu\nu}\Delta x^\mu\Delta x^\nu$ and the mass-shell value
$p_\mu p^\mu=m^2c^2$ are Lorentz invariants.

\begin{definition}[Covariance]\label{def:covariance}
A family of quantities $T^A$ is \emph{covariant} (transforms covariantly)
under $G$ if
\begin{equation}\label{covariant-object}
  g\cdot T^A = D(g)^A{}_B\,T^B
\end{equation}
for a representation $D$ of $G$. The equations \eqref{EOM} are
\emph{$\mathcal G$-covariant} if $E[\phi;b]$ is itself a covariant quantity
built from $(\phi,b)$, namely if there is a representation $D_E$ with
\begin{equation}\label{covariant-eq}
  E[\,g\cdot\phi;\;g\cdot b\,] \;=\; D_E(g)\,E[\phi;b]
  \qquad\text{for all } g\in\mathcal G \text{ and all } (\phi,b).
\end{equation}
\end{definition}

Equation~\eqref{covariant-eq} is the precise form of what older texts call
\emph{form invariance}: the same operator $E$, evaluated on the transformed
fields \emph{and} transformed background, is the transformed original. Note that
$g$ acts on $b$ as well; this is essential below.

\begin{definition}[Symmetry of the equations]\label{def:symmetry}
A transformation $g$ is a \emph{symmetry} of the system \eqref{EOM} at
fixed background $b_0$ if it maps solutions to solutions:
$E[\phi;b_0]=0\ \Rightarrow\ E[g\cdot\phi;\,b_0]=0$.
\end{definition}

\begin{proposition}[Covariance $+$ invariance of the background $=$ symmetry]
\label{prop:main}
Let the equations be $\mathcal G$-covariant in the sense of
\eqref{covariant-eq}, and let $b_0$ be a given background. Then the
stabilizer
\begin{equation}\label{stabilizer}
  G \;=\; \stab_{\mathcal G}(b_0)\;=\;\{\,g\in\mathcal G\;:\; g\cdot b_0=b_0\,\}
\end{equation}
is a symmetry group of $E[\phi;b_0]=0$.
\end{proposition}
\begin{proof}
For $g\in G$ and a solution $\phi$,
$E[g\cdot\phi;b_0]=E[g\cdot\phi;g\cdot b_0]=D_E(g)\,E[\phi;b_0]=0.$
\end{proof}

 Three different things
can be varied independently: the group $\mathcal G$ (how generous the
formulation is), the background $b_0$ (what is declared non-dynamical), and
the resulting physical symmetry group $G$.

\begin{remark}[Kretschmann's observation]\label{rem:kretschmann}
Covariance alone constrains nothing. Any equation can be made covariant under
\emph{any} group by promoting whatever breaks the covariance to a background
field and letting it transform along. Maxwell's equations written in
curvilinear coordinates, $\nabla_\mu F^{\mu\nu}=\tfrac{4\pi}{c}j^\nu$ with
the Christoffel symbols of the flat metric, are covariant under all
diffeomorphisms of $M$ provided $\eta_{\mu\nu}(x)$ is transformed as a
tensor field; yet the stabilizer of $\eta$ is still just the Poincar\'e
group. The physical content sits entirely in the stabilizer \eqref{stabilizer},
i.e.\ in what the background is and whether it is invariant. This was pointed
out by Kretschmann in 1917 as an objection to Einstein's claim that general
covariance had physical significance \cite{kretschmann,norton}; the resolution
is discussed in Sections~\ref{sec:gr} and~\ref{sec:kretschmann}.
\end{remark}

\begin{remark}[Larger covariance, smaller symmetry]\label{rem:smaller}
Enlarging $\mathcal G$ does not enlarge the symmetry group $G$. It only
enlarges the list of structures that must be declared as backgrounds, and $G$
is whatever preserves them. This is the ``absolute objects'' criterion of
Anderson \cite{anderson}: the symmetry group of a theory is the subgroup of
the covariance group that preserves its absolute (non-dynamical) objects.
\end{remark}

\begin{definition}[Variational symmetry]\label{def:variational}
Let the equations derive from an action, $E=\delta S/\delta\phi$. A
transformation $g$ is a \emph{symmetry of the action} (a variational
symmetry) if $S[g\cdot\phi;b_0]=S[\phi;b_0]$ up to boundary terms, i.e.\ if
the Lagrangian density changes by a total divergence.
\end{definition}

\begin{proposition}\label{prop:variational}
A variational symmetry is a symmetry of the equations. The converse is false.
\end{proposition}
\begin{proof}[Sketch]
Differentiate $S[g\cdot\phi]=S[\phi]$ in an arbitrary direction $\delta\phi$:
$\langle E[g\cdot\phi],\,Dg\,\delta\phi\rangle=\langle E[\phi],\delta\phi\rangle$,
where $Dg$ is the (invertible) linearization of $g$. Hence
$E[g\cdot\phi]=(Dg^{\dagger})^{-1}E[\phi]$, so $E[\phi]=0$ implies
$E[g\cdot\phi]=0$. For the converse, see Section~\ref{sec:summary} (rescaling
a linear equation; electromagnetic duality).
\end{proof}

Covariance of the \emph{action} means that $S$ is a scalar (the Lagrangian
is a scalar density) under $\mathcal G$ acting on $(\phi,b)$; together with
invariance of $b_0$ it yields a variational symmetry, which is the entry
point for Noether's theorem (Section~\ref{sec:noether}). Figure~\ref{fig:logic}
sketches the logical architecture.

\begin{figure}[ht]
\centering
\begin{tikzpicture}[
  box/.style={draw, rounded corners=2pt, align=center, inner sep=5pt,
              font=\small, text width=4.1cm},
  arr/.style={-{Stealth[length=2.4mm]}, thick}]
  \node[box] (cov) at (-3.6,0)    {Covariant equations\\ $E[\phi;b]=0$ under $\mathcal G$};
  \node[box] (inv) at ( 3.6,0)    {Invariant background\\ $g\cdot b_0=b_0$};
  \node[box] (sym) at (0,-2.4)    {Symmetry group\\ $G=\stab_{\mathcal G}(b_0)$};
  \node[box,text width=5cm] (n1)  at (-3.6,-4.9) {Global symmetry of the action\\ $\Rightarrow$ conserved currents\\
                                   (Noether I)};
  \node[box,text width=5cm] (n2)  at ( 3.6,-4.9) {Arbitrary functions of $x$ (no background)\\
                                   $\Rightarrow$ gauge redundancy, identities\\ (Noether II)};
  \draw[arr] (cov) -- (sym);
  \draw[arr] (inv) -- (sym);
  \draw[arr] (sym) -- (n1);
  \draw[arr] (sym) -- (n2);
\end{tikzpicture}
\caption{Logical structure of Proposition~\ref{prop:main} and its
consequences.}
\label{fig:logic}
\end{figure}


\section{Symmetry}
\label{sec:cm-sym}

Symmetry is often introduced to students as a single, monolithic idea: ``the
Lagrangian is invariant under such-and-such transformation, therefore there is
a conserved quantity, therefore the system respects the symmetry.'' In
practice, the situation is considerably richer, and the richness is not a
technicality but the source of some of the deepest phenomena in physics, from
the accidental degeneracies of the Kepler problem to the existence of the
electron mass and the pion.

It is useful to keep distinct objects in view:
\begin{enumerate}
\item the \textbf{Lagrangian} (or action, or Hamiltonian) that \emph{defines} 
    the dynamics;
\item the \textbf{equations of motion} (EOM) that it \emph{generates}; and
\item an actual \textbf{solution} of those equations, singled out by initial
    or boundary data;
\item the state that is actually realized as the ground state or equilibrium 
    configuration of the system.\footnote { Spontaneous Symmetry Breaking: 
        The Lagrangian, the equations of motion,
        \emph{and} the boundary conditions are all fully symmetric, and yet the
        stable, physically realized solution is not. This is not contingent: it
        is a structural consequence of the existence of a degenerate manifold
        of energy-minimizing (or free-energy-minimizing) configurations,
        together with an instability of the symmetric configuration and, in the
        case of a genuine phase transition, the thermodynamic (infinite-volume)
        limit.}

\end{enumerate}
Each of these objects has its own symmetry group, and the three groups need not
coincide. 

Let $q^a(t)$, $a=1,\dots,n$, be generalized coordinates on a 
configuration
space $Q$, with Lagrangian $L(q,\dot q,t)$ and action
\[
S[q] = \int_{t_1}^{t_2} L(q(t),\dot q(t),t)\,dt .
\]
\begin{equation*}
   p_a=\frac{ \partial L(q,\dot q,t)}{\partial \dot q_a}
\end{equation*}
The Euler--Lagrange (EL) operator acting on $L$ is
\[
[L]_a = \frac{\partial L}{\partial q^a} - 
\frac{d}{dt}\frac{\partial L}{\partial \dot q^a}
\]
and the equations of motion are $[L]_a=0$.

\begin{definition}[Variational / Noether symmetry]
An infinitesimal transformation $\delta q^a = \epsilon X^a(q,\dot q,t)$ is a
    \emph{variational symmetry} of $L$ if there exists a function $F(q,\dot
    q,t)$ such that, \emph{identically in $q(t)$} (not only on solutions),
    $\delta L=\epsilon \tfrac{dF}{dt}$. Equivalently, the action changes only by
    a boundary term, $\delta S = \epsilon\big[F\big]_{t_1}^{t_2}$, for
    arbitrary paths, not just extremal ones.
\end{definition}

Evidently adding the total derivative
\begin{equation*}
    \frac{d F (q,t)}{d t}=\frac{\partial F(q,t)}{\partial t} + \frac{ \partial F (q,t)}{\partial q_j} \dot q_j
\end{equation*}
to Lagrangian does not change equations of motion. However canonical momenta
become
\[ P_a=p_a+\frac{\partial F}{\partial q_a} \]
it is nothing but gauge invariance. Transformation $ p,q \rightarrow P,q $ is evidently
canonical and new Hamiltonian becomes 
\[ H_{new}=H(p(P),q)- \frac{\partial F}{\partial t}.\]

We are used to think that gauge invariance is about electromagnetism, actually it appears
already in classical mechanics. 
So canonical momenta have no physical meaning (since it is not gauge invariant).
This fact correspond to geometrical point of view: $p_a,q_a$ are merely local
coordinates on phase space.

Having EOM as system of ordinary differential equations,
we can transform it into equivalent one, say by multiplying every equation by arbitrary 
function which never becomes zero. This new system may correspond to another Lagrangian and
the difference between new and old one need not come down to total derivative. Most probably
new system will not be Lagrangian at all.

Let us look at two-dimensional harmonic oscillator.
Usual Lagrangian is
\begin{equation}
   \label {ho1}
    L_1= \frac{1}{2} (\dot x^2+\dot y^2) -\frac{ \omega^2}{2} ( x^2+ y^2)
\end{equation}
which yields equation  of motions:
\begin{equation}
   \label {hoe1}
   \begin{split}
       \ddot x+ \omega^2 x & =0 \\
       \ddot y+ \omega^2 y & =0 
   \end{split}
\end{equation}
Same equations of motion come form 
another Lagrangian
\begin{equation}
   \label {ho2}
    L_2= \frac{1}{2} (\dot x^2-\dot y^2) -\frac{ \omega^2}{2} ( x^2- y^2)
\end{equation}
or
\begin{equation}
   \label {ho3}
    L_3= \dot x\dot y - \omega^2  x y .
\end{equation}
Evidently these Lagrangians are not connected by adding total derivative.
Corresponding  Hamiltonians 
\begin{equation}
   \label {ho2h}
    H_2= \frac{1}{2} ( p_x^2-p_y^2) +\frac{ \omega^2}{2} ( x^2- y^2)
\end{equation}
and
\begin{equation}
   \label {ho3h}
    H_3= p_x p_y +\omega^2  x y
\end{equation}
are not even positive definite!

Two inclusions organize the entire subject:
\begin{itemize}
    \item $G_L \subseteq G_{EOM}$: \emph{every symmetry of the Lagrangian is a
        symmetry of the equations of motion, but the equations of motion can
        admit strictly more symmetry than is manifest in the Lagrangian one
        wrote down.} 
    \item $G_{sol} \subseteq G_{EOM}$: \emph{a particular solution, or a
        particular physically realized state, generally respects only a
        subgroup of the symmetry of the equations that produced it.} 
        Asymmetry is put in by hand
        through the choice of initial or boundary data. This is contingent:
        choosing symmetric data restores a symmetric solution, and averaging
        over the group orbit of solutions restores the full symmetry at the
        ensemble level.
   \end{itemize}

\begin{definition}[Symmetry of the equations of motion]
A transformation $g$ (finite or infinitesimal) is a \emph{symmetry of the EOM},
    or a \emph{dynamical symmetry}, if it maps every solution of $[L]_a=0$ to 
    another solution of the \emph{same} equations, i.e. it maps the solution 
    set to itself. No reference to the action or to boundary terms is required 
    by this definition.
\end{definition}

\begin{proposition}
\label{noether-forward}
Every variational symmetry of $L$ is a symmetry of the equations of motion. In
    symbols, $G_L \subseteq G_{EOM}$.
\end{proposition}

\begin{proof}[Proof (standard, physicist's level)]
Two facts are needed.
\emph{(a) The EL operator annihilates total time derivatives.} 
    identically  in $q(t)$ — not merely on shell.

\emph{(b) A change of path is a change of path, regardless of whether it is
    a symmetry.} Consider the one-parameter family of paths 
    $q_\epsilon(t) = q(t) + \epsilon X(q(t),\dot q(t),t) + O( \epsilon^2)$, 
    for \emph{arbitrary} $q(t)$, not necessarily a solution. By hypothesis,
\[
L(q_\epsilon,\dot q_\epsilon,t) = L(q,\dot q,t) + \epsilon\, \frac{dF}{dt} + O(\epsilon^2)
\]
identically in $q(t)$. Hence the action evaluated on $q_\epsilon$ differs from
    the action evaluated on $q$ only by the boundary term
    $\epsilon[F]_{t_1}^{t_2}$, \emph{for every path, not only extremal ones}.
    Consequently $q(t)$ is a stationary point of $S$ under variations that
    vanish at the endpoints if and only if $q_\epsilon(t)$ is — because the two
    actions differ by a term that does not depend on the interior of the path
    at all. Since ``being a solution of $[L]_a=0$'' is exactly the statement of
    being such a stationary point, $q(t)$ solves the EOM iff $q_\epsilon(t)$
    does, to first order in $\epsilon$. Thus the transformation maps solutions
    to solutions: it is a symmetry of the EOM.
\end{proof}

The familiar conserved charge follows from the same computation kept on shell:
\[
Q = \frac{\partial L}{\partial \dot q^a} X^a - F, 
\qquad \frac{dQ}{dt}\Big|_{\text{on shell}} = 0 .
\]

The converse of Proposition~\ref{noether-forward} fails in general: the
equations of motion routinely admit more symmetry than is manifest in the
Lagrangian one starts from. It is useful to separate \emph{manifest point
symmetries} of $L$ (transformations of $q$ and $t$ alone, geometrically visible
in the way $L$ is written), from \emph{generalized (dynamical) variational
symmetries}, which are allowed to depend on velocities or momenta and need not
correspond to any simple geometric transformation of configuration space, from
the full group $G_{EOM}$ of symmetries of the equations of motion:
\[
G_L^{\text{point}} \;\subseteq\; G_L^{\text{gen}} \;\subseteq\; G_{EOM} .
\]
It is a subtle question, belonging to the \emph{inverse problem of the calculus
of variations}, exactly when 
the second inclusion is an equality; under
favorable regularity assumptions many symmetries of a regular,
finite-dimensional Lagrangian system's EOM can indeed be exhibited as
\emph{generalized} variational symmetries (possibly of a Lagrangian differing
from the original by a boundary term), but this is not automatic, and 
in degenerate (constrained) systems it can fail outright. What
matters  is the \emph{first} inclusion: it is very
often strict, i.e., there is real, physically important symmetry of the EOM
that is simply invisible if one only stares at the naive, manifest symmetry of
$L$.

\begin{example}[Free particle: hidden $SL(2,\mathbf R)$]
The Lagrangian $L=\tfrac12 m\dot x^2$ manifestly has the Galilei symmetry
    (translations, boosts) plus time translation. But the equation of motion
    $\ddot x=0$ has straight lines as its general solution, and the space of
    straight lines is left invariant by the larger, one-dimensional-time
    conformal (``Schr\"odinger/M\"obius'') group acting projectively on $t$,
\[
t \;\longmapsto\; \frac{\alpha t+\beta}{\gamma t+\delta},
    \qquad \alpha\delta-\beta\gamma=1,
\]
together with a compensating rescaling of $x$. This $SL(2,\mathbf R)$ ``dynamical''
    symmetry (closely related to the conformal/Schr\"odinger symmetry of the
    free Schr\"odinger equation, and to the Ermakov--Lewis invariant of the
    time-dependent oscillator) is not visible as a point symmetry of
    $L=\tfrac12 m\dot x^2$ written in the ordinary way, yet it is a genuine
    symmetry of $\ddot x=0$.
\end{example}

\begin{example}[Kepler problem: hidden $SO(4)$]
For $L = \tfrac12 m\dot{\bm r}^2 + k/r$, the manifest symmetry is $SO(3)$ 
    (rotations) plus time translation. The equations of motion, however, 
    admit the additional conserved Laplace--Runge--Lenz vector
\[
\bm A = \bm p \times \bm L - mk\,\hat{\bm r},
\]
whose associated transformations close, together with angular momentum, into an
    $SO(4)$ algebra (for bound orbits). This hidden symmetry is responsible for
    the non-precession of Kepler ellipses and for the accidental $n^2$
    degeneracy of the hydrogen spectrum; it is a symmetry of the dynamics, not
    a geometrically manifest symmetry of the Lagrangian's potential term.
\end{example}

\begin{example}[Integrable systems: toroidal symmetry from action-angle variables]
More generally, for a Liouville-integrable Hamiltonian system with $n$ degrees
    of freedom, one can pass to action-angle variables $(I_i,\theta^i)$ in
    which $\mathcal{H} = \mathcal{H}(I)$ depends only on the actions. 
    The flow generated by
    $\mathcal{H}$ is then, trivially, an ($n$-torus of) simultaneous phase shifts
    $\theta^i \to \theta^i+c^i$, an abelian $U(1)^n$ symmetry group of the
    equations of motion. This symmetry is exact and completely general for
    integrable systems, yet it is essentially never visible as a point symmetry
    of the Lagrangian written in the original coordinates; the free-particle
    $SL(2,\mathbf R)$ and the Kepler $SO(4)$ above are, in this sense, special
    cases of the general toroidal symmetry structure of integrable systems,
    made visible only after a nontrivial (often nonlocal, momentum-dependent)
    change of variables.
\end{example}

\paragraph{The stabilizer of a solution}
Suppose $G_{EOM}$ acts on the solution space $\mathcal S$ of $[L]_a=0$.
 this action is well defined at least for
$g\in G_L$, and by definition for all of $G_{EOM}$). A single solution
$q(t)\in\mathcal S$, picked out by initial data $(q(t_0),\dot q(t_0))$ or by
boundary data, need not be a fixed point of this action. Its symmetry is the
\emph{stabilizer subgroup}
\[
    G_{sol}=\mathrm{Stab}_{G_{EOM}}(q) = \{ g\in G_{EOM} : g\cdot q
    = q \} \subseteq G_{EOM} .
\]
Generically $G_{sol} \subsetneq G_{EOM}$: the group moves the given
solution to a \emph{different} member of $\mathcal S$ (still a valid solution —
the equations do not care — but a different trajectory), rather than mapping it
to itself.

\begin{example}[Central force problem]
The Kepler/central-force EOM has full $SO(3)$ symmetry.
     A single
    orbit, however, has a definite, fixed angular momentum vector $\bm L = \bm
    r\times\bm p\neq 0$, determined by the initial position and velocity, and
    lies entirely in the plane orthogonal to $\bm L$. The stabilizer of this
    orbit inside $SO(3)$ is only the $SO(2)$ subgroup of rotations about the
    $\bm L$ axis (together with, for a closed orbit, possibly a discrete
    reflection). The other two rotational generators move the orbital plane to
    a different plane, i.e. map this solution to a genuinely different
    solution. Nothing is wrong with the dynamics; the reduction from $SO(3)$ to
    $SO(2)$ is entirely a consequence of the choice of initial data.
\end{example}

\begin{example}[Free particle]
$\ddot x=0$ is invariant under all spatial translations and Galilean boosts. A
    single trajectory $x(t)=x_0+v_0t$ is left invariant only by time
    translations (and, trivially, the identity among boosts/translations); a
    spatial translation or a boost maps it to a different straight line.
    Symmetric \emph{initial data} — e.g. $x_0=0,\,v_0=0$, the particle
    permanently at rest at the origin — is the unique choice whose stabilizer
    is the full translation-and-boost group; any other data breaks it
    explicitly.
\end{example}




\section{The infinitesimal form}
\label{sec:infinitesimal}

Let $\xi=\xi^\mu\partial_\mu$ generate a one-parameter subgroup $g_\eps$.
For a tensor field $T$ the active action to first order is
\begin{equation}\label{lie}
  g_\eps\cdot T = T-\eps\,\Lie_\xi T+O(\eps^2),
\end{equation}
where $\Lie_\xi$ is the Lie derivative; for a scalar,
$\delta\phi=-\eps\,\xi^\mu\partial_\mu\phi$, reproducing
\eqref{active}. The three notions become:
\begin{description}[leftmargin=1.2em,itemsep=3pt]
\item[Invariance of an object] $\Lie_\xi T=0$. For the background metric,
  $\Lie_\xi\eta_{\mu\nu}=\partial_\mu\xi_\nu+\partial_\nu\xi_\mu=0$ is the
  Killing equation. Its general solution
  $\xi^\mu=a^\mu+\omega^\mu{}_{\nu}x^\nu$, $\omega_{\mu\nu}=-\omega_{\nu\mu}$,
  has $4+6=10$ parameters: the Poincar\'e algebra.
\item[Covariance of an equation] $E[\phi;b]$ is a tensor field built
  \emph{naturally} from $\phi$ and $b$, so that
  \begin{equation}\label{naturality}
    \Lie_\xi E[\phi;b]=E'[\phi;b]\cdot\Lie_\xi\phi+\frac{\partial E}{\partial b}\cdot\Lie_\xi b,
  \end{equation}
  with $E'$ the linearization of $E$ in $\phi$.
\item[Symmetry] If $\Lie_\xi b_0=0$ and $\phi$ is a solution, then
  $E[\phi;b_0]=0$ implies $\Lie_\xi E=0$, and \eqref{naturality} reduces to
  \begin{equation}\label{zero-mode}
    E'[\phi;b_0]\cdot\Lie_\xi\phi=0.
  \end{equation}
  Thus $\Lie_\xi\phi$ solves the \emph{linearized} equations: a symmetry
  generates a one-parameter family of new solutions, $\phi_\eps=\exp(-\eps\Lie_\xi)\phi$.
\end{description}

\section{Symmetry of a law versus symmetry of a solution}
\label{sec:law-solution}

Definition~\ref{def:symmetry} concerns the \emph{law}. A separate and equally
important notion concerns an individual \emph{solution}.

\begin{definition}\label{def:isotropy}
The \emph{isotropy group} (or stabilizer) of a solution $\phi_0$ is
$G_{\phi_0}=\{g\in G: g\cdot\phi_0=\phi_0\}$, the subgroup of the symmetry group
under which \emph{that solution} is invariant. The set $G\cdot\phi_0
\cong G/G_{\phi_0}$ is the orbit of solutions generated from $\phi_0$.
\end{definition}

The distinction is the source of many apparent paradoxes:
\begin{itemize}[leftmargin=1.6em,itemsep=2pt]
\item The static Coulomb field of a point charge is invariant under rotations
  about the charge, $G_{\phi_0}=SO(3)\subset G$, although the law is invariant
  under the full Poincar\'e group. Translations generate the orbit of all
  Coulomb fields, one for each position of the charge.
\item A circular Kepler orbit in a central potential is invariant under
  neither time translations nor rotations about the orbital axis separately,
  but under the \emph{combined} one-parameter subgroup generated by
  $\partial_t+\omega\,\partial_\varphi$: a helical symmetry. Elliptic orbits
  have no continuous isotropy at all, although the law has $SO(3)\times
  \mathbb R_t$ (and more, Section~\ref{sec:kepler}).
\item In general relativity (Section~\ref{sec:gr}) the \emph{law} is invariant
  under all diffeomorphisms, but the \emph{solutions} typically have few or no
  Killing vectors; conserved quantities exist only for the symmetries of a
  given solution.
\item In spontaneous symmetry breaking the vacuum solution has a smaller
  isotropy group than the law (Section~\ref{sec:kepler}).
\end{itemize}
Symmetry of the law is a property of the theory; symmetry of the solution is a
property of a state (initial and boundary conditions included). Conservation
laws follow from the former, \emph{selection rules} and simplifying ans\"atze
from the latter.

\section{Noether's theorems}
\label{sec:noether}

\subsection{First theorem: global symmetries}

Consider a Lagrangian density $\mathcal L(\phi_a,\partial_\mu\phi_a)$ with
the background $b_0$ suppressed, and an infinitesimal transformation
$\delta\phi_a=\eps X_a$ under which
\begin{equation}\label{quasi}
  \delta\mathcal L=\eps\,\partial_\mu K^\mu
\end{equation}
(a variational symmetry in the sense of Definition~\ref{def:variational};
$K^\mu=0$ if $\mathcal L$ is strictly invariant). Writing $\delta\mathcal L$
in the standard way,
\begin{equation}
  \delta\mathcal L
  =\Big[\pd{\mathcal L}{\phi_a}-\partial_\mu\pd{\mathcal L}{(\partial_\mu\phi_a)}\Big]\delta\phi_a
   +\partial_\mu\Big(\pd{\mathcal L}{(\partial_\mu\phi_a)}\,\delta\phi_a\Big),
\end{equation}
and comparing with \eqref{quasi} gives, with $\mathrm{EL}_a$ the
Euler--Lagrange expressions,
\begin{equation}\label{noether1}
  \partial_\mu J^\mu=-\,\mathrm{EL}_a X_a\;\approx\;0,
  \qquad
  J^\mu=\pd{\mathcal L}{(\partial_\mu\phi_a)}X_a-K^\mu .
\end{equation}
Here $\approx$ means equality on solutions. The charge
$Q=\int J^0\,\dd^3x$ is conserved if the fields fall off at spatial infinity.

\begin{example}[Translations]\label{ex:translations}
A translation $x^\mu\to x^\mu+a^\mu$ acts as $\delta\phi_a=-a^\nu\partial_\nu\phi_a$.
If $\mathcal L$ has no explicit dependence on $x$ then
$\delta\mathcal L=-a^\nu\partial_\nu\mathcal L=\partial_\mu(-a^\mu\mathcal L)$,
so $K^\mu=-a^\mu\mathcal L$ and \eqref{noether1} gives the canonical
energy--momentum tensor,
\begin{equation}
  \partial_\mu T^\mu{}_\nu\approx0,\qquad
  T^\mu{}_\nu=\pd{\mathcal L}{(\partial_\mu\phi_a)}\,\partial_\nu\phi_a-\delta^\mu_\nu\,\mathcal L,
\end{equation}
with $T^0{}_0$ the energy density. If $\mathcal L$ depends explicitly on $x$,
which means precisely that a non-invariant background is present, then
$\partial_\mu T^\mu{}_\nu\approx-\,(\partial\mathcal L/\partial x^\nu)_{\rm expl}$
and the corresponding component is not conserved. Proposition~\ref{prop:main}
shows up here in the guise of Noether's theorem.
\end{example}

Table~\ref{tab:noether} lists the standard pairs of symmetry and conserved
quantity. Note that conservation requires a \emph{variational} symmetry; a
symmetry of the equations that is not variational produces no Noether charge
in this way.

\begin{table}[ht]
\centering
\small
\renewcommand{\arraystretch}{1.2}
\begin{tabular}{@{}lcl@{}}
\toprule
Symmetry & Parameters & Conserved quantity\\
\midrule
time translation & 1 & energy\\
space translations & 3 & momentum\\
rotations & 3 & angular momentum\\
Galilean boosts (quasi-invariance) & 3 & centre-of-mass motion, $m{\bm{x}}-{\bm{p}}\,t$\\
Lorentz boosts & 3 & boost charge $J^{0i}$\\
global phase $\phi\to e^{\ii\alpha}\phi$ & 1 & electric charge\\
dilatations (conformal theories) & 1 & dilatation charge\\
\bottomrule
\end{tabular}
\caption{Some symmetry--conservation pairs (Noether's first theorem).}
\label{tab:noether}
\end{table}

\subsection{Second theorem: local symmetries}

Suppose the action is invariant under transformations that depend on
\emph{arbitrary functions} $\eps^\alpha(x)$,
$\delta\phi_a=R_{a\alpha}\eps^\alpha$ with $R$ a differential operator. Then
$0=\delta S=\int \mathrm{EL}_a R_{a\alpha}\eps^\alpha\,\dd^4x$ for all
$\eps^\alpha$, and after integrating by parts
\begin{equation}\label{noether2}
  R^{\dagger}_{\alpha a}\,\mathrm{EL}_a\;\equiv\;0\qquad\text{(identically, off shell).}
\end{equation}
These \emph{Noether identities} are relations among the equations of motion
themselves, so the equations are not independent and part of the solution is
undetermined: gauge freedom.

\begin{example}[Electrodynamics]\label{ex:noether2-em}
Take
\begin{equation}\label{maxwell-action}
  S=\frac1c\int \dd^4x\Big[-\frac{1}{16\pi}F_{\mu\nu}F^{\mu\nu}-\frac1c\,j^\mu A_\mu\Big],
  \quad F_{\mu\nu}=\partial_\mu A_\nu-\partial_\nu A_\mu,
\end{equation}
with $\dd^4x=c\,\dd t\,\dd^3x$. Varying $A_\nu$ gives
$\partial_\mu F^{\mu\nu}=\tfrac{4\pi}{c}j^\nu$. The transformation
$\delta A_\mu=\partial_\mu\Lambda(x)$ leaves the field part invariant; the
Noether identity is $\partial_\nu\,\mathrm{EL}^\nu\equiv0$, i.e.\
$\partial_\nu\partial_\mu F^{\mu\nu}\equiv0$ on the field side. Consistency of
the field equations then demands $\partial_\nu j^\nu=0$: charge conservation
appears as a \emph{consistency condition}, not as the charge of a
global symmetry of the field.
\end{example}

\begin{remark}[Global symmetry versus gauge redundancy]\label{rem:gauge}
Global symmetries (finitely many constant parameters) map a configuration to a
\emph{physically different} one that is also a solution; they have conserved
charges that can be measured. Local transformations map a configuration to
a \emph{physically identical} one, written differently; they express
redundancy of the description, not a relation between distinct physical
situations. ``Gauge symmetry'' is thus a symmetry of the formulation in the
sense of covariance, and the word symmetry is used for it only by courtesy.
\end{remark}

\section{Example: Newtonian mechanics and the Galilei group}
\label{sec:galilei}

\subsection{Covariance and invariance of Newton's law}

Take $N$ particles with $m_a\ddot{{\bm{x}}}_a=\sum_{b\ne a}\bm F_{ab}$, where
$\bm F_{ab}=-\nabla_{a}V(|{\bm{x}}_a-{\bm{x}}_b|)$ is a central force. The Galilei group
(10 parameters) acts as
\begin{equation}\label{galilei}
  t'=t+s,\qquad {\bm{x}}'=R\,{\bm{x}}+{\bm{v}}\,t+\bm a,\qquad R\in SO(3).
\end{equation}
For \eqref{galilei},
\begin{equation}
  \ddot{{\bm{x}}}'_a=R\,\ddot{{\bm{x}}}_a,\qquad
  {\bm{x}}'_a-{\bm{x}}'_b=R\,({\bm{x}}_a-{\bm{x}}_b).
\end{equation}
The acceleration is thus \emph{invariant} under boosts and translations and
\emph{covariant} under rotations, and boosts and translations cancel in
differences of positions. Hence $\bm F'_{ab}=R\bm F_{ab}$ and the equation of
motion is covariant. The velocity $\dot{\bm{x}}'=R\dot{\bm{x}}+{\bm{v}}$, by contrast, is
\emph{not} covariant under boosts (it transforms inhomogeneously), so it can
enter a boost-symmetric force law only through differences of velocities.

\subsection{Breaking by a background: the medium}

Add friction against a medium with prescribed velocity field $\bm w({\bm{x}},t)$:
\begin{equation}\label{friction}
  m\ddot{\bm{x}}=-\gamma\,(\dot{\bm{x}}-\bm w).
\end{equation}
Under a boost, $\bm w\to\bm w'=\bm w+{\bm{v}}$, so \eqref{friction} is
\emph{covariant} under the Galilei group provided $\bm w$ is transformed
along with ${\bm{x}}$. But the boost is a symmetry only if the background is
unchanged. For a medium at rest ($\bm w=0$) the stabilizer consists of time
translations, space translations and rotations: boosts are excluded. The medium
plays the role of Anderson's absolute object, and its rest frame is a
preferred frame. Including the medium's atoms as dynamical degrees of
freedom restores the boost symmetry of the \emph{closed} system, exactly as in
Example~\ref{ex:coulomb}.

\subsection{Boosts: invariance up to a total derivative}

The Lagrangian $L=\sum\tfrac12m_a\dot{\bm{x}}_a^2-V$ is not invariant under boosts
but is \emph{quasi-invariant}: for the infinitesimal boost
$\delta{\bm{x}}_a=\eps\,\bm u\,t$,
\begin{equation}
  \delta L=\eps\sum_a m_a\bm u\cdot\dot{\bm{x}}_a=\eps\,\frac{\dd}{\dd t}\Big(\bm u\cdot\sum_a m_a{\bm{x}}_a\Big),
\end{equation}
which has the form \eqref{quasi} with $K=\bm u\cdot\sum m_a{\bm{x}}_a$. The
Noether charge $Q=\sum_a{\bm{p}}_a\!\cdot\delta{\bm{x}}_a/\eps-K=-\bm u\cdot\big(M\bm X-\bm P\,t\big)$
shows that
\begin{equation}
  \bm G=M\bm X_{\rm cm}-\bm P\,t=\text{const.}
\end{equation}
which says that the centre of mass moves uniformly. The quasi-invariance has a
deeper consequence: $\pb{G_i}{P_j}=M\delta_{ij}$, so the total mass $M$
appears as a \emph{central charge} of the Galilei algebra, in contrast to
the Poincar\'e algebra where $[K_i,P_j]\propto\delta_{ij}H$. This is why mass
is a conserved parameter of Galilean systems and why its quantum version
requires a projective representation (Section~\ref{sec:quantum}).

\subsection{Non-inertial frames and the geometric formulation}

Under a time-dependent change of frame ${\bm{x}}'=R(t){\bm{x}}+\bm a(t)$ Newton's law
acquires extra terms (Coriolis, centrifugal, Euler, and translational
inertial forces) and no longer has the same form. In the language of
Section~\ref{sec:definitions} the law is covariant under the Galilei group
\eqref{galilei} only. It can, however, be \emph{reformulated} so that it is
covariant under arbitrary coordinate changes: this is Cartan's
Newton--Cartan geometry \cite{cartan,dbkp}. Free fall is the geodesic equation
of a connection with $\Gamma^i{}_{00}=\partial_i\Phi$, and the absolute
objects are the time one-form $t_\mu=\partial_\mu t$ and the spatial
cometric $h^{\mu\nu}$. The formulation is generally covariant, the symmetry
is the subgroup preserving $t_\mu$ and $h^{\mu\nu}$, which is larger than the
Galilei group (it includes ${\bm{x}}\to{\bm{x}}+\bm a(t)$ acting on $\Phi$ by a term
linear in ${\bm{x}}$, proportional to $\ddot{\bm a}$, an expression of the
Newtonian equivalence principle). A textbook illustration of
Remark~\ref{rem:kretschmann}.

\section{Example: hidden symmetry and symmetry breaking}
\label{sec:kepler}

\subsection{Kepler problem: a symmetry that is not a coordinate transformation}

For $H={\bm{p}}^2/2m-k/r$ the geometric symmetry is $SO(3)\times\mathbb R_t$
(rotations and time translation), with conserved ${\bm{L}}={\bm{x}}\times{\bm{p}}$ and $H$.
There is, however, an additional conserved vector, the Runge--Lenz vector
\begin{equation}\label{rl}
  \bm A={\bm{p}}\times{\bm{L}}-mk\,\hat{{\bm{x}}},\qquad \pb{A_i}{H}=0 .
\end{equation}
The transformation it generates, $\delta{\bm{x}}=\eps\,\pb{{\bm{x}}}{{\bm{n}}\!\cdot\!\bm A}$,
\begin{equation}
  \delta{\bm{x}}=\eps\,\big[\,2({\bm{n}}\!\cdot\!{\bm{x}})\,{\bm{p}}-({\bm{x}}\!\cdot\!{\bm{p}})\,{\bm{n}}-({\bm{n}}\!\cdot\!{\bm{p}})\,{\bm{x}}\,\big],
\end{equation}
involves the \emph{momentum}: it is not a point transformation of
configuration space and is not related to any covariance of the
equation $m\ddot{\bm{x}}=-k{\bm{x}}/r^3$ under coordinate changes. It is nonetheless a
canonical symmetry of the Hamiltonian, and the Poisson algebra
\begin{equation}
  \pb{L_i}{L_j}=\epsilon_{ijk}L_k,\quad
  \pb{L_i}{A_j}=\epsilon_{ijk}A_k,\quad
  \pb{A_i}{A_j}=-2mH\,\epsilon_{ijk}L_k
\end{equation}
closes on $so(4)$ for $H<0$ (after the rescaling $\bm A\to\bm A/\sqrt{-2mH}$),
$e(3)$ for $H=0$, and $so(3,1)$ for $H>0$. The extra symmetry explains why
bound orbits are closed (the perihelion does not precess) \cite{goldstein}
and, in the quantum problem, the extra degeneracy of the hydrogen levels.
It is an example of a \emph{dynamical} (hidden) symmetry:
a symmetry of the law that is not visible as a covariance of the formulation.

\subsection{Spontaneous symmetry breaking}

An $O(N)$-invariant scalar theory,
$\mathcal L=\tfrac12\partial_\mu\phi_a\partial^\mu\phi_a-\tfrac\lambda4(\phi_a\phi_a-v^2)^2$,
has the symmetry group $O(N)$ \emph{for the law}. The lowest-energy
solutions form the sphere $\phi_a\phi_a=v^2$; each particular solution
$\phi_a=v\,\delta_{aN}$ has the smaller isotropy group $O(N-1)$. The orbit
$O(N)/O(N-1)=S^{N-1}$ is the vacuum manifold, along which fluctuations
cost no energy: $N-1$ massless Goldstone modes. The Noether currents remain
conserved; what changes is that the solution is not invariant. This is the
distinction of Section~\ref{sec:law-solution} once more.

\section{Common misunderstandings}

\begin{enumerate}[leftmargin=1.8em,itemsep=3pt,label=\textbf{\arabic*.}]
\item \textbf{``Covariant'' and ``invariant'' are not synonyms.} The fields
  ${\bm{E}},{\bm{B}}$ are covariant (they mix under boosts); $E^2-B^2$ and
  ${\bm{E}}\!\cdot\!{\bm{B}}$ are invariant. A covariant equation $\mathcal E=0$ is
  \emph{form-invariant}, while its left-hand side is not invariant but
  transforms in a representation. The standard phrase ``the equations are
  Lorentz invariant'' is acceptable shorthand only because it is clear from
  context that it is the \emph{solution set} that is invariant.
\item \textbf{Covariance is not physics; its stabilizer is.} Any equation can
  be written covariantly (Remark~\ref{rem:kretschmann}); what matters is
  which structures must then be declared as backgrounds, and whether they are
  invariant.
\item \textbf{A symmetry of the equations need not be a symmetry of the action.}
  The rescaling $x\to\lambda x$ of a \emph{linear} equation, such as the
  harmonic oscillator, maps solutions to solutions, but changes the action by
  $\lambda^2$; it has no Noether charge. Electromagnetic duality is another case.
  The converse (variational $\Rightarrow$ equation symmetry) holds
  (Proposition~\ref{prop:variational}).
\item \textbf{A symmetry of the law is not a symmetry of the solution.}
  Elliptic orbits, the Coulomb field, the vacuum of a broken theory, and the
  Schwarzschild metric all have smaller symmetry than their laws.
\item \textbf{Gauge ``symmetry'' is a redundancy.} Constant gauge
  transformations are symmetries with charges; local gauge transformations are
  not, and ``spontaneous breaking of a local symmetry'' is a statement about
  the \emph{description} and not about physical degeneracy.
\item \textbf{Transform everything.} An active symmetry transformation must act on
  the whole closed system (fields \emph{and} sources, apparatus,
  boundary conditions). Leaving a source untransformed amounts to introducing a
  background that breaks the symmetry (Example~\ref{ex:coulomb}).
\end{enumerate}

